\documentclass[12pt]{article}
\usepackage{bbm}
\usepackage{fancyhdr}
\usepackage{titling}
\usepackage{enumitem}
\usepackage{amsmath}
\usepackage{dsfont}
\usepackage{amsfonts}
\usepackage{amssymb}
\usepackage{amsthm}
\usepackage{geometry}
\usepackage{tikz}
\usepackage{algorithm}
\usepackage{algpseudocode}
\geometry{verbose,tmargin=1in,bmargin=1in,lmargin=1in,rmargin=1in}
\renewcommand{\vec}[1]{\mathbf{#1}}
\newcommand{\dec}{\text{dec}}
\newcommand{\divi}{\mathbin{\textnormal{div}}}
\title{Homework 1}
\author{Satvik Prasad}
\date{}
\newtheorem{lemma}{Lemma}[section]
\fancypagestyle{plain}{%
  \fancyhf{}
  \fancyhead[L]{\theauthor}
  \fancyhead[R]{\thedate}
}
\pagestyle{plain}
\makeatletter
\def\@maketitle{%
  \newpage
  \null
  \vskip 1em%
  \begin{center}%
    \let \footnote \thanks
    {\LARGE \@title \par}%
    \vskip 1em%
  \end{center}%
  \par
\vskip 1em}
\makeatother
\usepackage{cmbright}

\newcommand{\mat}[1]{\mathbf{#1}}

% ---------- Probability theory ----------
% Core operators (optional argument-free; use as \E{X}, \Var{X}, \Cov{X}{Y})
\DeclareMathOperator{\EOp}{\mathbb{E}}
\DeclareMathOperator{\VarOp}{Var}
\DeclareMathOperator{\CovOp}{Cov}
\DeclareMathOperator{\CorrOp}{Corr}

\renewcommand{\P}[1]{\mathbb{P}\left(#1\right)}        % \P{A}
\newcommand{\E}[1]{\EOp\left[#1\right]}                % \E{X}
\newcommand{\Var}[1]{\VarOp\left(#1\right)}            % \Var{X}
\newcommand{\Cov}[2]{\CovOp\left(#1,\,#2\right)}       % \Cov{X}{Y}
\newcommand{\Corr}[2]{\CorrOp\left(#1,\,#2\right)}     % \Corr{X}{Y}
\newcommand{\SD}[1]{\operatorname{SD}\left(#1\right)}  % \SD{X}

% Subscripted versions, e.g. \Esub{\theta}{X}, \Psub{\theta}{A}
\newcommand{\Esub}[2]{\EOp_{#1}\left[#2\right]}
\newcommand{\Psub}[2]{\mathbb{P}_{#1}\left(#2\right)}

% Conditioning, indicators, independence
\newcommand{\given}{\,\middle|\,}                      % \E{X \given Y}
\newcommand{\ind}[1]{\mathbbm{1}\left\{#1\right\}}     % \ind{X > 0}
\newcommand{\indep}{\perp\!\!\!\perp}                  % X \indep Y
\newcommand{\iid}{\overset{\text{iid}}{\sim}}          % X_i \iid \Normal(0,1)

% Convergence
\newcommand{\convp}{\xrightarrow{\;p\;}}               % in probability
\newcommand{\convd}{\xrightarrow{\;d\;}}               % in distribution
\newcommand{\convas}{\xrightarrow{\text{a.s.}}}        % almost surely
\newcommand{\convLp}[1]{\xrightarrow{L^{#1}}}          % in L^p

% Probability spaces
\newcommand{\F}{\mathcal{F}}                           % sigma-algebra
\newcommand{\B}{\mathcal{B}}                           % Borel sigma-algebra
\newcommand{\pspace}{(\Omega, \F, \mathbb{P})}

% Common distributions
\newcommand{\Normal}{\mathcal{N}}
\newcommand{\Bern}{\operatorname{Bernoulli}}
\newcommand{\Bin}{\operatorname{Binomial}}
\newcommand{\Pois}{\operatorname{Poisson}}
\newcommand{\Unif}{\operatorname{Uniform}}
\newcommand{\Exp}{\operatorname{Exp}}
\newcommand{\Geom}{\operatorname{Geometric}}
\newcommand{\Gam}{\operatorname{Gamma}}
\newcommand{\Beta}{\operatorname{Beta}}

% Number sets
\newcommand{\R}{\mathbb{R}}
\newcommand{\N}{\mathbb{N}}
\newcommand{\Z}{\mathbb{Z}}


\begin{document}
\begin{enumerate}
    \item For the following use inner product $$\langle f, g \rangle = \int_{0}^1 f(x) g(x) \, dx.$$ We want to find a polynomial $p \in \mathcal P_3(\mathbb R)$ such that $p(0) = p'(0)$ and minimizes
    $$
    \int_{0}^{1} (2 + 3x - p(x))^2 \, dx.
    $$
    This is an inner product over $\mathcal P_3(\mathbb R),$ meaning we want to minimize $$
    \langle 2 + 3x - p(x), 2 + 3x - p(x) \rangle = ||2 + 3x - p(x)||^2,
    $$
    given $p(0) = p'(0).$ Letting $p(x) = a+bx + cx^2 + dx^3,$ we have $a = b,$ from $p(0) =p'(0),$ meaning $$
    p(x) = a + ax + cx^2 + dx^3
    $$
    for scalars $a, c, d\in \mathbb R.$ Therefore, we look to minimize
    $$
    ||(2 - a) + (3 - a)x - cx^2 - dx^3||^2 = ||2 + 3x  - (a + ax + cx^2 + dx^3)||^2.
    $$

    We seek to project $2 + 3x$ onto the subspace $U = \{a + ax + cx^2 + dx^3 : a, c, d \in \mathbb R\}.$ We know the projection $P_U$ satisfies
    $$
    ||v - P_U v||^2 \leq  ||v - u||^2
    $$
    for all $u \in U,$ meaning setting $p = P_U (2 + 3x)$  will minimize the objective.

    Applying Graham-Schmidt to the basis of $U$ given by $\{1 + x, x^2, x^3\}$ yields the basis 
    $$
    \{ \sqrt{3/7} \cdot (1 + x), \sqrt{15/13} \cdot (4x^2 - x  - 1), \sqrt{3/7} \cdot (35x^3 -35 x^2 + 2x + 2)\},
    $$
    yielding decomposition for $2 + 3x$ given by $$
    P_U(2+3x) = \langle 2+ 3x, e_1 \rangle e_1 + \langle 2 + 3x, e_2\rangle e_2 + \langle 2 +3x, e_3\rangle e_3,
    $$
    resolving to 
    $$
    \begin{aligned}
    P_U(2+3x) &= \frac{11}{2}\cdot\frac{3}{7}\,(1+x) + \frac{1}{6}\cdot\frac{15}{13}\,(4x^2 - x - 1) - \frac{1}{12}\cdot\frac{3}{7}\,(35x^3 - 35x^2 + 2x + 2) \\
    &= \frac{33}{14}(1+x) + \frac{5}{26}(4x^2 - x - 1) - \frac{1}{28}(35x^3 - 35x^2 + 2x + 2) \\
    &= \frac{381}{182} + \frac{381}{182}x + \frac{105}{52}x^2 - \frac{5}{4}x^3
    \end{aligned}
    $$
    as required.

    \item Suppose that $T \in \mathcal L(V)$ and that $U$ is a subspace of $V.$ 
    
    First, let us assume that $U$ is invariant under $T,$ meaning for all $u \in U,$ $Tu \in U.$ Then, consider $T^\star w$ for some $w \in U^{\perp}.$ By the definition of the adjoint, we know for all $u \in U,$
    $$
    \langle u, T^\star w\rangle = \langle Tu, w \rangle.
    $$
    Since $Tu \in U$ and $w \in U^{\perp},$ this is equal to $0,$ meaning
    $$
    \langle u, T^\star w\rangle = 0,
    $$
    and hence $T^\star w \in U^\perp,$ allowing us to conclude $T^\star$ is invariant under $U^\perp.$

    Conversely, assume that $T^\star$ is invariant under $U^\perp.$ Then $T^\star w \in U^{\perp}$ for all $w \in U^{\perp}$. Since we know for all $u \in U,$ $w \in U^{\perp}$ $$
    \langle Tu, w\rangle = \langle u, T^\star w \rangle = 0,
    $$
    we can conclude that $Tu \in (U^\perp)^\perp = U$ since $w \in U^\perp,$ meaning $Tu \in U,$ i.e., $T$ is invariant under $U.$

    Therefore, $T$ is invariant under $U$ iff $T^\star$ is invariant under $U^\perp.$
\end{enumerate}
\end{document}